\documentclass[10pt,a4paper]{amsart}
\usepackage[margin=19mm]{geometry}
\usepackage{amsmath,amssymb,mathtools}
\usepackage[scr=rsfs,cal=euler]{mathalfa}
\usepackage{xfrac}
\usepackage[unicode,hidelinks,bookmarks=false]{hyperref}
\newcommand{\ie}{i.e.\ }
\newcommand{\cf}{cf.\ }
\newcommand{\resp}{respectively\ }
\newcommand{\Subsection}{\subsection}
\providecommand{\nobreakdash}{\nobreak\hbox{-}}
\setlength{\emergencystretch}{2em}
\multlinegap0pt
\usepackage{amssymb}
\usepackage{float}
\usepackage{cancel}
\usepackage{mathtools}
\newtheorem{lemma}{Lemma}
\newcommand{\dd}{\mathrm{d}}
\title{A series of Nash resolutions of a singular foliation}
\author{Ruben Louis}
\date{Source: arXiv:2301.08706v2 (CC BY 4.0)}
\begin{document}
\maketitle

\noindent\emph{Source and licence.} A series of Nash resolutions of a singular foliation. Version arXiv:2301.08706v2 (CC BY 4.0), distributed by arXiv under the Creative Commons Attribution 4.0 International licence, http://creativecommons.org/licenses/by/4.0/, as recorded in the arXiv licence metadata for this exact version and shown on the abstract page https://arxiv.org/abs/2301.08706v2. Full licence notice: \texttt{licenses/arxiv-2301.08706-CC-BY-4.0.txt}.\par\smallskip

\noindent\textit{Editorial scope.} Counted unit: the article's lemma lemma:dimensions (source0.tex lines 877--879). The article's complete original proof of it is delivered with it (source0.tex lines 880--898, one \texttt{proof} environment of 19 lines). No mathematics is written, repaired or re-derived: the statement and the proof are the article's own lines, in the article's own notation, and no retained line is substituted or re-flowed.  No further source-local context is needed: every cross-reference of every retained line resolves inside the two delivered spans.  Nothing was omitted from any retained span, so the package carries no omission record.  No retained line contains a citation, so the wrapper carries no bibliography and no bibliography entry.  No retained line contains a figure environment, an image-inclusion command or drawing code, so no image is delivered and the package declares no asset dependency.  Editorial formatting only, in addition: an AMS \texttt{amsart} preamble that restates the article's own declarations and macros used by the retained lines, namely the environments \texttt{lemma} on separate counters, together with the standard packages those lines need.  The retained environments print in this document as Lemma 1 in order of appearance, whereas the article prints the same items with the numbering of its own full document.  The complete native source of the article as distributed in the arXiv e-print for this exact version is bundled unchanged as \texttt{sources/136346/source0.tex}.
\begin{lemma}\label{lemma:dimensions}
{Let $(E, \dd, \rho)$ and $(E', \dd', \rho')$ be geometric resolutions of $\mathfrak F$. For all $i\geq 1$,  and for all  $V\in(\pi_i^{E})^{-1}(x)$ and $V'\in(\pi_i^{E'})^{-1}(x)$, one has  $\dim \overline{V}=\dim \overline{V'}$.}
\end{lemma}

\begin{proof}
If $x\in M$ is a regular point, then $\overline{V}=\overline{V'}=\{0\}$. Thus, the equality holds. Let $x\in M$ be a singular point. We prove it only for  $i=1, 2$, since  $i=1$ is a special case and  for $i\geq 3$ the proof uses a similar argument as for the one of $i=2$. The key point in the latter is, for every $x\in M$, the restriction of the complexes $(E, \dd, \rho)$ and $(E', \dd', \rho')$ at $x$ are quasi-isomorphic. This implies that  the codimension of $\mathrm{im}\left(\dd_x^{(i+1)}\right)$ inside $\ker \dd_x^{(i)}$, resp. $\mathrm{im}\left({\dd_x'}^{(i+1)}\right)$ inside $\ker {\dd_x'}^{(i)}$, is invariant.

Let $V\in(\pi_1^{E})^{-1}(x)$ and $V'\in(\pi_1^{E'})^{-1}(x)$. We have
\begin{align*}
  \dim \overline{V}&=\dim V-\dim (\mathrm{im}\,(\dd_x^{(2)}))\\&=\dim V-(\dim \ker \rho_x-\dim\ker \rho'_x+\dim (\mathrm{im}\,({\dd_x'}^{(2)}))\\&= \dim V -\mathrm{rk}({E_{-1}})+\mathrm{rk}({E'_{-1}})-\dim (\mathrm{im}\,({\dd_x'}^{(2)}))\\&=\dim V'-\dim (\mathrm{im}\,({\dd'_x}^{(2)}))\\&=\dim\overline{V'}.
\end{align*}
We have used the fact the cohomology groups at degree $-1$ of both complexes are isomorphic and the Rank–nullity theorem.

For $i=2$, let $V\in(\pi_2^{E})^{-1}(x)$ and $V'\in(\pi_2^{E'})^{-1}(x)$. Notice that $\dim V=\mathrm{rk}(E_{-2})-\mathrm{rk}(E_{-1})+r$. We have a similar formula for $\dim V'$. By direct computation we find that
\begin{align}
 \nonumber\dim \overline{V}&= \dim V- \dim (\mathrm{im}\,\dd_x^{(3)})\\\label{eq:invariant}&=\dim V-\mathrm{rk}(E_{-2})+ \mathrm{rk}(E'_{-2})+ \dim (\mathrm{im}\,(\dd_x^{(2)}))- \dim (\mathrm{im}\,({\dd_x'}^{(2)}))- \dim (\mathrm{im}\,({\dd_x'}^{(3)})).
\end{align}
We have used the fact the cohomology groups at degree $-2$ of both complexes are isomorphic and the Rank–nullity theorem. But $$\dim (\mathrm{im}\,(\dd_x^{(2)}))=\mathrm{rk}(E_{-1})-\dim (\mathrm{im}(\rho_x))- \dim W,$$ where $W$ is such that $\dim (\mathrm{im}\,(\dd_x^{(2)}))\oplus W=\ker \rho_x$. A similar formula holds for $\dim (\mathrm{im}\,({\dd'_x}^{(2)}))$ by adding ${'}$ everywhere. Substituting them into the Equation \eqref{eq:invariant}  we obtain

$$\dim \overline{V}= \dim \overline {V'}+ \dim W'-\dim W= \dim \overline {V'},$$ since $\dim W'=\dim W$.


\end{proof}

\end{document}
